\documentclass{article}
\usepackage{amsmath}
\usepackage{amssymb}
\usepackage{amsthm}
\usepackage{booktabs}
\usepackage{tikz}
\usepackage{lua-ul}
\usepackage[a4paper,margin=25mm]{geometry}

\newtheorem{theorem}{Theorem}[section]
\newcommand{\term}[1]{\textbf{#1}}
\newcommand{\engine}{Fermion TeX Engine}

\begin{document}

\tableofcontents

\section{Live LuaLaTeX}
\label{sec:live}

This document is typeset by a \emph{real} Lua\LaTeX{} --- genuine
Knuth--Plass line breaking, real Latin Modern fonts with ligatures and
kerning, and the complete \term{amsmath} machinery.\footnote{This footnote is typeset live: the engine captures
TeX's insert nodes and the resident page builder reserves space for them
at the bottom of the page, rule and all.} The \engine{} keeps
the document resident: when you edit one sentence, only the block
containing it is re-typeset by LuaTeX, and only the pages whose display
list changed are patched in the preview.

Edit any word in this paragraph and watch the inspector: one block
compiles, every other block is reused from cache, and untouched pages
survive by reference. Efficient office workflows -- ligatures intact.
\underLine{This sentence is underlined by lua-ul and should appear in
the live preview.}

\section{Mathematics}
\label{sec:math}

Inline mathematics like $\int_0^\infty e^{-x^2}\,dx = \frac{\sqrt{\pi}}{2}$
is typeset by TeX itself, inside the paragraph. Numbered displays work
with genuine equation counters:

\begin{equation}
  \label{eq:gauss}
  \nabla \cdot \mathbf{E} = \frac{\rho}{\varepsilon_0}
\end{equation}

\begin{align}
  \label{eq:euler}
  e^{i\pi} + 1 &= 0, \\
  \sum_{n=1}^{\infty} \frac{1}{n^2} &= \frac{\pi^2}{6}.
\end{align}

Equation~\eqref{eq:gauss} is Gauss's law; the pair in~\eqref{eq:euler}
needs \term{align}. Insert a new numbered equation above them and the
numbers below renumber through the counter chain --- the engine
recompiles exactly the blocks whose entry state changed.

\begin{theorem}
  \label{thm:main}
  A resident typesetting runtime patches a page if and only if the
  page's display list hash changes.
\end{theorem}

Theorem~\ref{thm:main} lives in Section~\ref{sec:math}; its number
comes from a real \term{amsthm} counter injected per block.

\section{Graphics and Tables}
\label{sec:graphics}

Figure~\ref{fig:plot} is a genuine float: declared here, placed by the
live page builder according to its \texttt{[t]} specifier, with a real
caption counter. Edit its coordinates and only this block recompiles:

\begin{figure}[t]
\centering
\begin{tikzpicture}[scale=1.1]
  \draw[->, thick] (-0.2,0) -- (3.4,0) node[right] {$x$};
  \draw[->, thick] (0,-0.2) -- (0,2.2) node[above] {$y$};
  \draw[blue, very thick, domain=0:1.4, smooth] plot (\x, {\x*\x}) node[right] {$y = x^2$};
  \draw[red, very thick, domain=0:2.1, smooth, dashed] plot (\x, {sqrt(\x)}) node[right] {$y = \sqrt{x}$};
  \fill[purple] (1,1) circle (2.2pt) node[above left] {$(1,1)$};
\end{tikzpicture}
\caption{A live float: TikZ inside \texttt{figure}, placed at the top of
the page by the resident output routine.}
\label{fig:plot}
\end{figure}

Tables use \term{booktabs} rules, exactly as lualatex would print them:

\begin{center}
\begin{tabular}{lrrr}
\toprule
Layer & Nodes & Compiled & Reused \\
\midrule
Source DOM   & 24 & 1 & 23 \\
Chunk cache  & 24 & 1 & 23 \\
Page DOM     & 3  & 1 & 2  \\
\bottomrule
\end{tabular}
\end{center}

\section{Incrementality}
\label{sec:inc}

As stated in Section~\ref{sec:live}, the engine is a runtime, not a
compiler wrapper. Blocks are diffed by content hash; each dirty block
is typeset in isolation with its entry state (section, equation and
theorem counters) injected, and its exit state chains into the next
block. Labels are captured from the real aux stream, so
\eqref{eq:gauss} and \ref{thm:main} resolve exactly as in a full
compile.

Renaming a label makes precisely the referencing blocks dirty. A
whitespace-only edit re-typesets one block, produces a chunk that is
byte-identical to the cached one, and therefore patches nothing. The
inspector proves both claims on every keystroke.

Structural changes --- a new package in the preamble, a class option
--- take the honest path: the preamble format is rebuilt and every
block recompiles once. Everything else stays live.

\section{References in Motion}
\label{sec:refs}

Citations resolve from the live bibliography: \cite{knuth84} defined
the machine, and \cite{lamport94} made it a document system. Renaming
or reordering the entries below renumbers these citations exactly like
a rerun of latex would.

\begin{thebibliography}{9}
\bibitem{knuth84} Donald E. Knuth. \emph{The \TeX book}.
Addison--Wesley, 1984.
\bibitem{lamport94} Leslie Lamport. \emph{\LaTeX: A Document
Preparation System}. Addison--Wesley, 1994.
\end{thebibliography}

\end{document}
